% Copyright (c) 2026 Marcio Luiz Lima de Oliveira and V. el Hakim.
% Distribution and modification are subject to LPPL version 1.3c.
% SPDX-License-Identifier: LPPL-1.3c
% Part of the LPPL-covered work listed in MANIFEST.txt.
% See LICENSE and LICENSING.txt for the licence and scope.

\chapter{Figures, tables and algorithms}\label{chap:03}
\renewcommand{\FigLocation}{Chapters/Chapter_3/fig/}
\renewcommand{\ChapLocation}{Chapters/Chapter_3/}
\begin{absbox}
Keep assets beside their chapter, cite their origin and show enough information for readers to understand them. These examples demonstrate photographs, subfigures, tables, diagrams and pseudocode.
\end{absbox}

\section{A figure stored with its chapter}
Every chapter sets \verb|\FigLocation| before its figures. Use a filename with
an explicit extension and a width relative to the current line, especially
inside a narrow box or subfigure. \autoref{fig:ch3:earth} uses a downloaded NASA
Apollo 17 photograph, stored in this chapter's \texttt{fig} directory.

\begin{figure}[htbp]
\centering
\includegraphics[width=0.52\linewidth]{\FigLocation earth-apollo17.jpg}
\caption[Earth from Apollo 17]{Earth from Apollo 17. Credit: NASA/Apollo 17 crew,
7 December 1972; public-domain image obtained through Wikimedia Commons.
Source and download details are in \texttt{fig/ATTRIBUTION.txt}.}
\label{fig:ch3:earth}
\end{figure}

\begin{lstlisting}
\renewcommand{\FigLocation}{Chapters/Chapter_3/fig/}
\begin{figure}[htbp]
  \centering
  \includegraphics[width=0.52\linewidth]
    {\FigLocation earth-apollo17.jpg}
  \caption[Short list entry]{A descriptive caption and credit.}
  \label{fig:ch3:earth}
\end{figure}
\end{lstlisting}
Place \verb|\label| after \verb|\caption| so it records the figure number.
The optional caption supplies a shorter list-of-figures entry. The placement
letters permit a figure here, at the top, at the bottom, or on a float page.
They are preferences, not fixed positions. Use \verb|\clearpage| at a deliberate
boundary to flush waiting floats instead of repeatedly forcing placement.

\section{Compare panels without creating a new asset}
The existing layout loads \texttt{subfig}. Use \verb|\subfloat| from that package;
do not also load \texttt{subcaption}. The second panel below crops the same image
at typesetting time. This is a layout demonstration, not a new observation.
\begin{figure}[htbp]
\centering
\subfloat[Full photograph]{\includegraphics[width=0.39\linewidth]{\FigLocation earth-apollo17.jpg}}
\hfill
\subfloat[Central crop]{\includegraphics[width=0.39\linewidth,trim=30mm 30mm 30mm 30mm,clip]{\FigLocation earth-apollo17.jpg}}
\caption{A full image and a crop of the same NASA/Apollo 17 photograph. Credit
and public-domain source as in \autoref{fig:ch3:earth}.}
\label{fig:ch3:panels}
\end{figure}
\begin{lstlisting}
\subfloat[Full photograph]{\includegraphics[width=.39\linewidth]
  {\FigLocation earth-apollo17.jpg}}
\hfill
\subfloat[Central crop]{\includegraphics[width=.39\linewidth,
  trim=30mm 30mm 30mm 30mm,clip]{\FigLocation earth-apollo17.jpg}}
\end{lstlisting}
The trim order is left, bottom, right, top; unitless values are big points.
This example uses explicit millimetres. Keep the total crop smaller than the
image's natural dimensions, which depend on its resolution metadata.
Always state when a crop, contrast adjustment or annotation changes how an image
is presented. Keep the original asset and its attribution with the chapter.

\section{Tables with readable precision}
\autoref{tab:ch3:results} contains invented data, solely to demonstrate a table.
\texttt{booktabs} supplies the horizontal rules; \texttt{siunitx} supplies the
\texttt{S} column that aligns decimal values. Put units in column headings.
\begin{table}[htbp]
\centering
\caption{Illustrative measurements; these are not research results.}
\label{tab:ch3:results}
\begin{tabular}{@{}l S[table-format=2.1] S[table-format=1.2]@{}}
\toprule
Method & {Time (s)} & {Error (a.u.)}\\\midrule
Baseline & 12.4 & 0.32\\
Alternative & 10.8 & 0.28\\
Reference & 15.0 & 0.25\\\bottomrule
\end{tabular}
\end{table}
\begin{lstlisting}
\begin{tabular}{@{}l S[table-format=2.1]@{}}
  \toprule
  Method & {Time (s)} \\
  \midrule
  Baseline & 12.4 \\
  \bottomrule
\end{tabular}
\end{lstlisting}
Use \texttt{tabularx} for wrapping text columns, \texttt{longtable} for a table
that can break over pages, and \texttt{multirow} only when a shared row heading
actually helps. A \texttt{longtable} is not placed inside a \texttt{table} float.
Avoid shrinking a wide table until the letters become unreadable.

\section{Editable diagrams}
\begin{figure}[htbp]
\centering
\begin{tikzpicture}[node distance=8mm,>=stealth]
\node[draw=\ColorTheme,rounded corners,minimum height=9mm] (input) {Input};
\node[draw=\ColorTheme,rounded corners,minimum height=9mm,right=of input] (method) {Method};
\node[draw=\ColorTheme,rounded corners,minimum height=9mm,right=of method] (output) {Output};
\draw[->,thick] (input) -- (method);
\draw[->,thick] (method) -- (output);
\end{tikzpicture}
\caption{A small vector diagram drawn with TikZ and the current chapter colour.}
\label{fig:ch3:workflow}
\end{figure}
For a diagram this small, its TikZ source can remain in the chapter. Larger
diagram sources can live in the same \texttt{fig} directory and be included
with \verb|\input{\FigLocation workflow.tex}|. PDF is suitable for exported
vector artwork; PNG and JPEG are suitable for raster images. No EPS conversion
or shell escape is needed for the examples in this template.

\section{Pseudocode and ordinary footnotes}
\begin{algorithm}[htbp]
\caption{Mean of a nonempty sequence}\label{alg:mean}
\begin{algorithmic}[1]
\Require Values $u_1,\ldots,u_n$ with $n>0$
\Ensure Arithmetic mean $m$
\State $s\gets 0$
\For{$i=1,\ldots,n$}
  \State $s\gets s+u_i$
\EndFor
\LineComment{Normalise only after accumulating the observations.}
\State $m\gets s/n$
\State \Return $m$
\end{algorithmic}
\end{algorithm}
\texttt{algorithm} supplies the floating caption; \texttt{algorithmic} supplies
the pseudocode. The configuration adds \verb|\LineComment{text}| as a full-line
comment. Refer to it with \verb|\autoref{alg:mean}|, as in \autoref{alg:mean}.
Use \verb|\myfootnote{text}| for an ordinary numbered footnote.\myfootnote{This
is a normal footnote; its helper has the same argument as \texttt{\string\footnote}.}
Publication-origin notes use a separate helper introduced in \autoref{chap:05}.
